\chapter{Canonical Quantization Examples}
\label{cqe}

\section {Canonical Quantization}
\label{canonical-quantization}
Let start from some definitions \footnote{
 There is some confusion in terminology, exact  operator domain 
 rarely important in physical problems, so physicists generally do 
 not pay attention to difference between Hermitian and self-adjoint operators.}:
\begin{definition}
    Operator $ \hat A^\dag$ is called adjoint to operator $ \hat A $,
    if 
$ \langle \hat A\phi,\psi \rangle = \langle \phi,\hat A^\dag \psi \rangle$.
\end{definition}
Note that domain of $ \hat A^\dag$ includes that of $ \hat A $
$ D( \hat A)\subseteq D( \hat A^\dag)$.
\begin{definition}
    Operator $ \hat A $ is called Hermitian if $ \hat A = \hat A^\dag $ on
    $D( \hat A)$.
\end{definition}
\begin{definition}
    Operator $ \hat A $ is called  self-adjoint if $ \hat A = \hat A^\dag $
    and $D( \hat A) = D( \hat A^\dag)$.
\end{definition}
For bounded operators on Hilbert space $ \mathcal{H} $ 
$D( \hat A) = \mathcal{H} = D( \hat A^\dag)$ so Hermitian
and self-adjoint bounded operators are the same. 

Canonical quantization assigns Hermitian (self-adjoint) operators
$\hat q_i, \hat p_j$ on Hilbert space $ \mathcal{H} $
satisfying canonical commutation relations
\begin{equation}
    \label{cqe:ccr}
    [\hat q_i, \hat p_j] = i\hbar \delta_{ij}.
\end{equation}
to classical observables $q_i, p_j$ --- coordinates on phase space $ \mathbb{R}^{2 n} $.
All other observables are obtained by substitution of these operators into
classical observables --- functions on phase space.
Generally operator ordering ambiguities arise,  
and additional analytic input is required to ensure self-adjointness.
This prescription is insufficient for general systems, it works only
if system is described by global Cartesian coordinates.
So the standard name "canonical quantization" is misleading, but generally used.

Operator $ \hat H$ (Hamiltonian) generates time evolution (unitary group):
    \begin{equation}
       \label {cqe:evol}
 \hat U(t) = e^{-i \hat H t/ \hbar}.
    \end{equation}
    Evolution may be seen as action on observables (Heisenberg picture)
\begin{align}
   \hat A(t)&= \hat U(t) \hat A \hat U^{-1}(t), \label {cqe:heis1}\\
    \frac{d \hat A}{d t} &= \frac{i}{\hbar}[\hat H,\hat A],\label{cqe:heis2}
\end{align}
    or as action on states (Schr\"odinger picture)
\begin{align}
    |\psi(t)\rangle&= \hat U^{-1}(t) |\psi \rangle, \label{cqe:schr1} \\
    \frac{d |\psi \rangle}{d t} &=-\frac{i}{\hbar} \hat H |\psi \rangle. \label{cqe:schr2}
\end{align}

For "usual" Hamiltonian function $ H=p^2/(2 m) +V(q) $ we get:
\begin{equation}
   \label {cqe:ham}
    \hat H = \frac{\hat p^2}{2 m}+V(\hat q)
\end{equation}

In particular using \eqref{cqe:ham},\eqref{cqe:ccr} and \eqref{cqe:heis2} we get:
\begin{equation}
   \label {cqe:ehrenfest}
   \begin{split}
       \frac{d \hat q_i}{d t} &= \frac{\hat p_i}{m}\\
       \frac{d \hat p_j}{d t} &= - \frac{\partial V}{\partial q_j}(\hat q)
   \end{split}
\end{equation}
which formally coincide with classical equation of motion.

For most practical problems canonical quantization is realized as
\begin{equation}
q_i \mapsto \hat q_i = q_i, \qquad p_j \mapsto \hat p_j = -i\hbar \frac{\partial}{\partial q_j},
\end{equation}
acting on $\mathcal{H}=L^2(Q,dq)$. \footnote{ Actually this representation
is essentially unique.}

Below we shall consider a number of very simple but instructive 
examples.

\section{Quantum Dynamics of a Wave Packet in a Square Well}

Consider a particle confined to a one-dimensional infinite square well of width $L$,
\begin{equation}
V(x) = 
\begin{cases}
0, & 0 < x < L,\\
\infty, & \text{otherwise}.
\end{cases}
\end{equation}

A \emph{wave packet} is a superposition of stationary states,
\begin{equation}
\Psi(x,0) = \sum_{n=1}^{\infty} c_n \, \phi_n(x), 
\qquad
\phi_n(x) = \sqrt{\frac{2}{L}} \sin \frac{n \pi x}{L},
\end{equation}
where the coefficients $c_n$ encode the initial localization and momentum distribution of the particle.  
Time evolution follows from the Schr\"odinger equation,
\begin{equation}
\Psi(x,t) = \sum_{n=1}^{\infty} c_n \, \phi_n(x) \, e^{-i E_n t/\hbar}, 
\qquad
E_n = \frac{n^2 \pi^2 \hbar^2}{2 m L^2}.
\end{equation}
The quadratic energy spectrum is responsible for the characteristic quantum features described below.

For a wave packet narrowly localized in both position and momentum (centered around a large $n_0$):
\begin{itemize}
    \item The packet moves approximately according to the classical trajectory,
    \begin{equation}
    x_{\mathrm{cl}}(t) = x_0 + v t,
    \end{equation}
    reflecting elastically at the walls.
    \item Ehrenfest's theorem ensures that the expectation values of position and momentum follow classical equations of motion at short times.
\end{itemize}

Due to the quadratic spectrum $E_n \sim n^2$, different stationary components acquire phases at different rates.  
Consequences include:
\begin{itemize}
    \item Spreading (dispersion) of the wave packet.
    \item Destruction of the initial localization by interference.
    \item A characteristic \emph{collapse time} after which the probability density appears irregular.
\end{itemize}
The evolution remains fully unitary; no energy loss or dissipation occurs.

The exact quadratic spectrum leads to wave packet \emph{revivals}:
\begin{equation}
T_{\mathrm{rev}} = \frac{4 m L^2}{\pi \hbar}.
\end{equation}
At $t = T_{\mathrm{rev}}$, the wave packet reconstructs itself up to a global phase,
\begin{equation}
\Psi(x,T_{\mathrm{rev}}) = e^{i \theta} \Psi(x,0),
\end{equation}
while at rational fractions $t = \frac{p}{q} T_{\mathrm{rev}}$, \emph{fractional revivals} occur, splitting the packet into $q$ localized copies.

Inside the infinite well:
\begin{itemize}
    \item A self-adjoint momentum operator compatible with the boundary conditions does not exist.
    \item Operationally, momentum must be described via a POVM, not a PVM.
    \item The packet behaves as if momentum reverses sign at the walls, producing strong interference near the boundaries.
\end{itemize}

The dynamics of a wave packet in a square well illustrate:
\begin{description}
    \item [Short-time classical motion:] The packet moves and reflects approximately like a classical particle.
    \item [Quantum dispersion:] Interference between stationary components spreads the packet.
    \item [Exact revivals:] The wave packet periodically reconstructs itself due to the quadratic spectrum.
    \item [Unitary boundary-induced interference:] Reflections at the walls generate characteristic oscillations without dissipation.
\end{description}
This system exemplifies the interplay between classical intuition and fundamentally quantum effects.

